\chapter{The art of API design}
\label{ch:apidesign}

\begin{aside}
A library's API is its UI for programmers. Same principles
apply: minimise surface, name well, default sensibly.
\end{aside}

\section{Minimise surface}

Every public symbol is a maintenance burden. Hide
implementation; expose intent.

\section{Name well}

Names communicate purpose. \code{set\_visible(true)} is
clearer than \code{set\_v(1)}. Verb-noun for actions.

\section{Sensible defaults}

The default should be what 80\% of users want. Make the
common case zero-config.

\section{Builder vs constructor}

For elements with many optional parameters, the pipe
operator builder is friendlier than a 10-arg constructor:

\begin{lstlisting}
text("hi") | bold | fg<"red"> | width(20);
\end{lstlisting}

Versus:

\begin{lstlisting}
text("hi", true, "red", 20, ...);
\end{lstlisting}

\section{Errors as values}

Use \code{std::expected} or a sentinel return; avoid
exceptions for routine failures.

\section{Composability}

Functions that take and return elements compose. Functions
that take config structs less so.

\section{Versioning}

Public API changes break users. Use semver; deprecate
gradually.

\section{Recap}

\begin{recap}
\begin{itemize}[leftmargin=*]
  \item Small surface, clear names.
  \item Defaults for the common case.
  \item Builders for many optional args.
  \item Errors as values; not exceptions for routine.
  \item Compose, don't accumulate.
\end{itemize}
\end{recap}

\section*{Workshop}

\begin{enumerate}[leftmargin=*]
  \item Audit your own public API for over-exposure.
  \item Rename one badly-named function.
  \item Convert one config struct to a builder.
\end{enumerate}
